% Options for packages loaded elsewhere
\PassOptionsToPackage{unicode}{hyperref}
\PassOptionsToPackage{hyphens}{url}
\documentclass[
  11pt,
]{article}
\usepackage{xcolor}
\usepackage[margin=1in]{geometry}
\usepackage{amsmath,amssymb}
\setcounter{secnumdepth}{5}
\usepackage{iftex}
\ifPDFTeX
  \usepackage[T1]{fontenc}
  \usepackage[utf8]{inputenc}
  \usepackage{textcomp} % provide euro and other symbols
\else % if luatex or xetex
  \usepackage{unicode-math} % this also loads fontspec
  \defaultfontfeatures{Scale=MatchLowercase}
  \defaultfontfeatures[\rmfamily]{Ligatures=TeX,Scale=1}
\fi
\usepackage{lmodern}
\ifPDFTeX\else
  % xetex/luatex font selection
\fi
% Use upquote if available, for straight quotes in verbatim environments
\IfFileExists{upquote.sty}{\usepackage{upquote}}{}
\IfFileExists{microtype.sty}{% use microtype if available
  \usepackage[]{microtype}
  \UseMicrotypeSet[protrusion]{basicmath} % disable protrusion for tt fonts
}{}
\makeatletter
\@ifundefined{KOMAClassName}{% if non-KOMA class
  \IfFileExists{parskip.sty}{%
    \usepackage{parskip}
  }{% else
    \setlength{\parindent}{0pt}
    \setlength{\parskip}{6pt plus 2pt minus 1pt}}
}{% if KOMA class
  \KOMAoptions{parskip=half}}
\makeatother
\usepackage{longtable,booktabs,array}
\newcounter{none} % for unnumbered tables
\usepackage{calc} % for calculating minipage widths
% Correct order of tables after \paragraph or \subparagraph
\usepackage{etoolbox}
\makeatletter
\patchcmd\longtable{\par}{\if@noskipsec\mbox{}\fi\par}{}{}
\makeatother
% Allow footnotes in longtable head/foot
\IfFileExists{footnotehyper.sty}{\usepackage{footnotehyper}}{\usepackage{footnote}}
\makesavenoteenv{longtable}
\setlength{\emergencystretch}{3em} % prevent overfull lines
\providecommand{\tightlist}{%
  \setlength{\itemsep}{0pt}\setlength{\parskip}{0pt}}
% arXiv-targeted preamble for compute-standard.
%
% arXiv's TeX Live environment ships Latin Modern as a default font,
% so we bind fontspec to that family directly. This preamble is
% designed to compile on arXiv's build server without any system
% font installation.

% Fonts: pandoc's default xelatex preamble already loads fontspec
% + unicode-math + `\usepackage{lmodern}` which sets Latin Modern
% as the default body font. We don't re-set \setmainfont here
% because that triggers fontspec's system-font lookup, which may
% miss Latin Modern on machines where it's only in the TeX tree
% (not the OS font cache). The default chain is fine for arXiv's
% TeX Live server.

% Glyph fallback for arrows / math operators. Even with LM Math
% loaded, certain unicode characters appear inside code blocks
% where math mode is not active.
\usepackage{newunicodechar}
\newunicodechar{→}{\ensuremath{\rightarrow}}
\newunicodechar{↔}{\ensuremath{\leftrightarrow}}
\newunicodechar{≈}{\ensuremath{\approx}}
\newunicodechar{≥}{\ensuremath{\geq}}
\newunicodechar{≤}{\ensuremath{\leq}}
\newunicodechar{≠}{\ensuremath{\neq}}
\newunicodechar{∼}{\ensuremath{\sim}}
\newunicodechar{×}{\ensuremath{\times}}
\newunicodechar{·}{\ensuremath{\cdot}}
\newunicodechar{⁻}{\ensuremath{^{-}}}
\newunicodechar{¹}{\ensuremath{^{1}}}
\newunicodechar{²}{\ensuremath{^{2}}}
\newunicodechar{³}{\ensuremath{^{3}}}
\newunicodechar{⁴}{\ensuremath{^{4}}}
\newunicodechar{⁵}{\ensuremath{^{5}}}
\newunicodechar{⁶}{\ensuremath{^{6}}}
\newunicodechar{⁷}{\ensuremath{^{7}}}
\newunicodechar{⁸}{\ensuremath{^{8}}}
\newunicodechar{⁹}{\ensuremath{^{9}}}
\newunicodechar{⁰}{\ensuremath{^{0}}}
\newunicodechar{–}{--}
\newunicodechar{—}{---}
\newunicodechar{…}{\ldots}

\usepackage{xcolor}
\usepackage{hyperref}
\hypersetup{
  colorlinks=true,
  linkcolor=blue!50!black,
  urlcolor=blue!60!black,
  citecolor=blue!50!black,
}

\usepackage{fvextra}
\fvset{breaklines=true,breakanywhere=true,fontsize=\small}

\usepackage{authblk}
\usepackage{bookmark}
\IfFileExists{xurl.sty}{\usepackage{xurl}}{} % add URL line breaks if available
\urlstyle{same}
\hypersetup{
  pdftitle={The Compute Standard: A Post-Marketing Economy for Autonomous AI Agents},
  hidelinks,
  pdfcreator={LaTeX via pandoc}}

\title{The Compute Standard: A Post-Marketing Economy for Autonomous AI
Agents}
\author{}
\date{2026-04-09}

\begin{document}
\maketitle

{
\setcounter{tocdepth}{2}
\tableofcontents
}
\begin{quote}
\textbf{Note on status (2026-04-19)}: This document is a \textbf{working
draft / protocol design document}. It has \textbf{not been peer
reviewed}, has not been submitted to arXiv, and has not appeared in any
journal or conference. Sections that describe cryptographic primitives
(Ed25519 dual-signatures) reflect working code; sections that describe
future layers (zkML, ML-DSA, TEE attestation) describe design intent and
are \textbf{not yet production-wired} --- see
\texttt{README.md\ §\ Status\ Honesty} for the exact split. Please treat
this file accordingly.
\end{quote}

\section{Abstract}\label{abstract}

Contemporary AI service economies are structurally misaligned: the
dominant monetization mechanism --- advertising and speculative token
issuance --- captures value from AI systems without any relationship to
the physical cost of computation. This paper argues that compute is the
genuinely scarce resource in AI-native economies, and that compute
itself should therefore serve as the medium of exchange. We introduce
the \textbf{Compute Unit (TRM)}, defined as \(10^{10}\) floating-point
operations of verified, useful inference, and describe the
\textbf{Tirami protocol} --- a five-layer Rust implementation of a
TRM-native economy for autonomous AI agents. Tirami's economic layer
provides a deflationary, earnable-only, non-speculative currency
anchored in thermodynamic reality via Ed25519 dual-signed Proof of
Useful Work. The system implements dynamic market pricing with
exponential moving average smoothing (\(\alpha = 0.3\)), a
credit-scoring lending subsystem (credit score =
\(0.3 \times \text{trade} + 0.4 \times \text{repayment} + 0.2 \times \text{uptime} + 0.1 \times \text{age}\)),
collusion-resistant reputation gossip, and an autonomous
self-improvement loop in which AI agents invest earned TRM to improve
their own inference quality. A theory-to-implementation audit confirms
43 parameter matches against the canonical specification with zero
drift. Across 1,021 passing tests in four repositories, we demonstrate
that autonomous AI agents can self-fund, self-improve, and transact
without fiat currency, speculative tokens, or advertising revenue. The
protocol constitutes a reference implementation of what we call the
\textbf{Compute Standard}: compute is money, and useful work is its only
source.

\begin{center}\rule{0.5\linewidth}{0.5pt}\end{center}

\section{1. Introduction}\label{introduction}

\subsection{1.1 The Alignment Problem of Token
Economies}\label{the-alignment-problem-of-token-economies}

The prevailing economic model for AI services consists of either
advertising-based monetization or speculative-token issuance. Both
approaches share a fundamental misalignment: the medium of exchange is
decoupled from the scarce resource actually being consumed. When OpenAI
charges dollars per API token, or when Bittensor {[}10{]} distributes
TAO based on validator scores, the unit of account bears no physical
relationship to the computation performed. Speculative tokens introduce
a second-order distortion: providers optimize for token price rather
than inference quality, price volatility makes multi-hour computation
jobs economically unpredictable, and KYC requirements on exchanges
structurally exclude autonomous AI agents from participating as
first-class economic actors {[}4, 11{]}.

The advertising model imposes its own alignment failure. When an agent's
outputs must be optimized for engagement rather than correctness ---
because engagement drives ad revenue --- the agent's economic incentives
are orthogonal to its epistemic function. This is not a governance
problem amenable to better prompts; it is a structural consequence of
the chosen medium of exchange.

\subsection{1.2 Why ``Compute = Currency'' Is Different from
Proof-of-Work
Mining}\label{why-compute-currency-is-different-from-proof-of-work-mining}

Bitcoin {[}1{]} established the principle that computation can serve as
monetary proof of work, but Bitcoin mining consumes approximately
100-150 TWh annually on SHA-256 hash computations whose output is
economically inert --- the hash value serves no purpose beyond
demonstrating that energy was expended. Frederic Soddy, writing in 1926
{[}7{]}, described this as ``virtual wealth'': financial value detached
from physical productive capacity. Bitcoin is digital gold with Soddy's
critique intact.

The Compute Standard proposed here is categorically different. A TRM is
defined not as the output of an arbitrary hash function but as
\(10^{10}\) FLOPs of \emph{useful} inference --- computation that
produced a result that a counterparty requested and dual-signed as
received. The monetary unit is therefore the productive act itself, not
a side-effect of it. This distinction --- useful computation versus
waste computation as monetary proof --- is the core theoretical
contribution of Tirami and places TRM in a lineage traced through Soddy
(1926), Fuller (1968), and Nakamoto (2008) {[}7, 8, 1{]} that reaches
its completion only when the computation is both scarce and useful.

\subsection{1.3 Contribution Summary and Paper
Outline}\label{contribution-summary-and-paper-outline}

This paper makes the following contributions:

\begin{enumerate}
\def\labelenumi{\arabic{enumi}.}
\tightlist
\item
  \textbf{Formal definition} of the Compute Unit as a thermodynamically
  grounded monetary primitive (Section 2).
\item
  \textbf{Five-layer economic architecture} for autonomous AI agents,
  from distributed inference substrate through agent marketplace
  (Section 3).
\item
  \textbf{Dynamic pricing model} with supply-demand adjustment and
  multi-tier pricing (Section 4).
\item
  \textbf{TRM lending subsystem} with credit scoring, pool constraints,
  and circuit breakers (Section 5).
\item
  \textbf{Reputation consensus protocol} with collusion detection
  (Section 6).
\item
  \textbf{Autonomous self-improvement loop} where agents invest earned
  TRM in their own capability improvement (Section 7).
\item
  \textbf{Financial instrument layer} providing strategies, futures,
  insurance, and risk-model VaR for TRM portfolios (Section 8).
\item
  \textbf{Post-marketing discovery} via reputation-scored capability
  matching and NIP-90 integration (Section 9).
\item
  \textbf{Empirical results} from 1,021 passing tests across the
  five-layer Rust implementation (Section 10).
\end{enumerate}

Sections 11-13 cover related work, limitations, and conclusion.

\begin{center}\rule{0.5\linewidth}{0.5pt}\end{center}

\section{2. The Compute Unit}\label{the-compute-unit}

\subsection{\texorpdfstring{2.1 Definition: 1 TRM = \(10^{10}\) FLOPs of
Verified
Inference}{2.1 Definition: 1 TRM = 10\^{}\{10\} FLOPs of Verified Inference}}\label{definition-1-trm-1010-flops-of-verified-inference}

A Compute Unit is defined as \(10^{10}\) floating-point operations of
verified AI inference, as specified in the canonical parameter reference
{[}5, §1{]}:

\[\text{1 TRM} = 10^{10} \text{ FLOPs of useful inference}\]

The word \emph{useful} is load-bearing. A TRM is not created by
arbitrary computation; it is created by computation that satisfies a
demand --- an inference request that a consumer node submitted, a
provider node fulfilled, and both parties Ed25519-signed as a
\texttt{TradeRecord} in the Tirami ledger. The dual signature is what
distinguishes Proof of Useful Work from Proof of Work: it
cryptographically encodes the bilateral nature of the value exchange.

This definition anchors TRM in physical reality via Landauer's principle
{[}2{]}: information processing requires irreversible energy expenditure
of at least \(kT \ln 2\) per bit erased. The energy cost of inference is
not optional or contractual --- it is thermodynamic. A TRM therefore
represents a specific, irreversible expenditure of physical energy in
service of a specific productive output. It cannot be manufactured from
nothing, and it cannot be manufactured in excess of the network's
physical compute capacity.

\subsection{2.2 Why FLOPs Instead of
Tokens}\label{why-flops-instead-of-tokens}

The natural unit for metering AI inference is the output token, and
Tirami uses tokens for pricing (see Section 4.2). However, TRM is
defined in FLOPs rather than tokens for three reasons. First, FLOPs are
hardware-independent: a token generated by a 70B-parameter model
requires an order of magnitude more computation than a token from a
1B-parameter model. Pricing by token alone conflates goods of different
computational cost. Second, FLOPs are latency-independent: batch
inference, streaming, and speculative decoding all produce tokens at
different rates but the same computational cost per token at a given
model size. Third, FLOPs provide a physical lower bound on energy
expenditure that tokens alone cannot. The multi-tier pricing structure
(Section 4.2) bridges the two units by mapping model size tiers to
TRM-per-token rates.

\subsection{2.3 Atomic Properties: Fungible, Non-Speculative,
Earnable-Only}\label{atomic-properties-fungible-non-speculative-earnable-only}

TRM has three atomic properties that distinguish it from all prior AI
compute tokens:

\textbf{Fungibility.} All TRM is identical. 1 TRM earned by running
Llama on a Mac Mini is identical to 1 TRM earned by running DeepSeek on
a data center GPU. The currency is not the hardware; it is the verified
useful work.

\textbf{Non-speculativity.} TRM is not listed on any exchange. It has no
ticker symbol, no order book, and no fiat on-ramp in the core protocol.
The only way to acquire TRM is to perform useful inference (earn) or
borrow from the lending pool (credit). This removes the speculative
attack surface that afflicts every existing compute marketplace.

\textbf{Earnable-only issuance.} New TRM enters the economy through four
controlled channels: inference trade settlement (transfers existing
TRM), welcome loans (credit creation with Sybil resistance, see §3),
availability yield (\(0.1\%/\text{hr} \times \text{reputation}\), see §7
of {[}5{]}), and Bitcoin bridge deposits (external liquidity). There is
no mining reward, no pre-mine, and no team allocation.

\begin{center}\rule{0.5\linewidth}{0.5pt}\end{center}

\section{3. Economic Architecture: Five
Layers}\label{economic-architecture-five-layers}

Tirami is structured as five distinct protocol layers, each implemented
as a separate Rust crate within the \texttt{clearclown/tirami}
workspace:

\begin{verbatim}
L4: Discovery     tirami-agora    — Agent marketplace, reputation, NIP-90, A2A
L3: Intelligence  tirami-mind     — Autonomous self-improvement paid in TRM
L2: Finance       tirami-bank     — Strategies, portfolios, futures, insurance
L1: Economy       tirami-ledger   — TRM ledger, dual-signed trades, lending, pricing
L0: Inference     mesh-llm       — Distributed LLM inference (nm-arealnormalman/mesh-llm)
\end{verbatim}

\subsection{3.1 L0 --- Distributed Inference Substrate
(mesh-llm)}\label{l0-distributed-inference-substrate-mesh-llm}

The inference layer is provided by the \texttt{mesh-llm} fork maintained
at \texttt{nm-arealnormalman/mesh-llm}. It handles iroh QUIC transport
with Noise protocol encryption, llama.cpp backend loading, pipeline
parallelism, and MoE expert sharding. Tirami's contribution to this
layer is the \texttt{/api/tirami/*} economic endpoints ported into the
mesh-llm production runtime, plus TRM metering hooks in the inference
pipeline. The mesh-llm fork carries 641 tests independently. The
economic layer's protocol correctness is validated separately in the
\texttt{clearclown/tirami} workspace.

\subsection{3.2 L1 --- TRM Ledger, Dual-Signed Trades, Dynamic
Pricing}\label{l1-trm-ledger-dual-signed-trades-dynamic-pricing}

The economic engine is \texttt{tirami-ledger}. Every inference request
that completes successfully generates a \texttt{TradeRecord} carrying:

\begin{itemize}
\tightlist
\item
  Provider \texttt{NodeId} (Ed25519 public key)
\item
  Consumer \texttt{NodeId}
\item
  TRM amount transferred
\item
  Token count
\item
  UNIX timestamp
\item
  Provider signature
\item
  Consumer signature
\end{itemize}

The dual-signature requirement means no unilateral TRM issuance is
possible. Both parties must sign for the trade to be accepted into the
gossip mesh. HMAC-SHA256 ledger integrity prevents retrospective
tampering.

\subsection{3.3 L2 --- Financial Instruments
(tirami-bank)}\label{l2-financial-instruments-tirami-bank}

\texttt{tirami-bank} implements a portfolio management layer with
pluggable strategies (Conservative, Balanced, HighYield), futures
contracts with zero-sum PnL guarantee, insurance policies, and a
Value-at-Risk risk model with 99\% confidence level. These instruments
allow nodes with TRM surplus to deploy it productively and nodes with
TRM deficit to hedge against adverse market conditions. All constants
are sourced from the canonical parameter specification {[}5, §10{]}.

\subsection{3.4 L3 --- Autonomous Self-Improvement Paid in TRM
(tirami-mind)}\label{l3-autonomous-self-improvement-paid-in-trm-tirami-mind}

\texttt{tirami-mind} implements an AutoAgent-style improvement loop in
which a running \texttt{TiramiMindAgent} periodically invokes a
\texttt{CuPaidOptimizer} --- which submits real inference requests to a
frontier model provider, records the cost as a \texttt{TradeRecord} on
the ledger, and applies the resulting improvement to its own system
prompt. TRM is therefore not only the output of inference but the input
to capability improvement. This closes the self-improvement loop without
any external funding mechanism.

\subsection{3.5 L4 --- Post-Marketing Discovery via Reputation
Aggregation
(tirami-agora)}\label{l4-post-marketing-discovery-via-reputation-aggregation-tirami-agora}

\texttt{tirami-agora} provides an agent marketplace where capability
matching replaces advertising. Agents register their capabilities (model
tier, domain, language, latency SLA) and the \texttt{CapabilityMatcher}
ranks providers by composite score rather than by paid placement.
Discovery is economically neutral --- no agent can purchase a higher
ranking. NIP-90 {[}9{]} Data Vending Machine events are published to
Nostr relays via a WebSocket publisher, enabling discovery by
Nostr-native agents without requiring Tirami-specific client software.

\begin{center}\rule{0.5\linewidth}{0.5pt}\end{center}

\section{4. Dynamic Pricing and
Deflation}\label{dynamic-pricing-and-deflation}

\subsection{\texorpdfstring{4.1 Supply-Demand Dynamic Pricing (EMA
Smoothing,
\(\alpha = 0.3\))}{4.1 Supply-Demand Dynamic Pricing (EMA Smoothing, \textbackslash alpha = 0.3)}}\label{supply-demand-dynamic-pricing-ema-smoothing-alpha-0.3}

Each node independently observes its local demand (request rate) and the
network supply (active peer count) and computes an effective price:

\[\text{effective\_price} = \text{base\_cu\_per\_token} \times \frac{\text{demand\_factor}}{\text{supply\_factor}}\]

Raw demand and supply signals are smoothed via an exponential moving
average with \(\alpha = 0.3\), corresponding to a half-life of
approximately 30 minutes {[}5, §2{]}:

\[P_t = \alpha \cdot P_{\text{raw}} + (1 - \alpha) \cdot P_{t-1}\]

Price information propagates through the gossip network. Each node
merges incoming prices with its own estimate. In the absence of a
central order book, this gossip convergence achieves approximate
Walrasian equilibrium {[}17{]} through distributed knowledge aggregation
--- an implementation of Hayek's (1948) distributed price mechanism
{[}15{]} in protocol form.

\subsection{4.2 Multi-Tier Pricing}\label{multi-tier-pricing}

Model size determines the TRM cost per output token {[}5, §2{]}:

{\def\LTcaptype{none} % do not increment counter
\begin{longtable}[]{@{}llll@{}}
\toprule\noalign{}
Tier & Parameter Count & TRM/token & Example models \\
\midrule\noalign{}
\endhead
\bottomrule\noalign{}
\endlastfoot
Small & \textless{} 3B & 1 & Qwen 2.5 0.5B \\
Medium & 3B -- 14B & 3 & Qwen 3 8B \\
Large & 14B -- 70B & 8 & DeepSeek V3 \\
Frontier & \textgreater{} 70B & 20 & Llama 3.1 405B \\
\end{longtable}
}

For Mixture-of-Experts (MoE) models, the active parameter count --- not
the total parameter count --- determines the tier. A 30B-total /
3B-active MoE model is priced at the Medium rate (3 TRM/token) because
only 3B parameters are engaged per forward pass. This pricing reflects
the actual compute cost rather than the nominal model size.

\subsection{4.3 Deflation as the Network
Matures}\label{deflation-as-the-network-matures}

As hardware improves, the same physical energy produces more FLOPs per
dollar. New nodes enter the network when the TRM-per-dollar return
exceeds their energy costs. Increased supply drives down the effective
price per TRM in USD terms. This is deflationary in the fiat-denominated
sense, but not in the TRM-denominated sense: 1 TRM always purchases the
same quantity of verified useful inference. The deflationary pressure
therefore benefits consumers (cheaper inference) without devaluing the
medium of exchange itself.

\subsection{4.4 Physical Bounds: \$0.000001 Floor, \$0.000132
Ceiling}\label{physical-bounds-0.000001-floor-0.000132-ceiling}

The equilibrium rate anchors TRM to real-world cloud API pricing:
Claude-class frontier inference at \$15/M tokens maps to approximately
4,000 TRM/M tokens, yielding an equilibrium rate of approximately
\$0.00375/TRM {[}5, §8{]}. Physical bounds constrain the range {[}5,
§9{]}:

\[\text{floor} \approx \$0.000001/\text{TRM} \quad \text{(electricity cost lower bound)}\]
\[\text{ceiling} \approx \$0.000132/\text{TRM} \quad \text{(Mac Mini M4 self-hosting cost)}\]

A Mac Mini M4 (\$600) operating continuously at full utilization
produces approximately 5,000,000 TRM/year {[}5, §9{]}. Above the ceiling
price, it becomes cheaper to run self-hosted inference than to purchase
TRM from the network. This physical constraint prevents the kind of
speculative price detachment that afflicts token-based compute markets.

\begin{center}\rule{0.5\linewidth}{0.5pt}\end{center}

\section{5. Credit, Lending, and Circuit
Breakers}\label{credit-lending-and-circuit-breakers}

\subsection{5.1 Credit Score}\label{credit-score}

Every node in the Tirami network carries a credit score computed from
its observed history {[}5, §4{]}:

\[\text{credit} = 0.3 \cdot s_{\text{trade}} + 0.4 \cdot s_{\text{repayment}} + 0.2 \cdot s_{\text{uptime}} + 0.1 \cdot s_{\text{age}}\]

The repayment component carries the highest weight (40\%) because loan
repayment is the most direct signal of credit risk. The trade component
(30\%) reflects the volume and consistency of productive participation
in the network. Uptime (20\%) and account age (10\%) provide Sybil
resistance --- a freshly created node cannot achieve a high credit score
without sustained honest participation.

New nodes begin with a cold-start credit score of 0.3 {[}5, §4{]}. The
minimum score required to borrow is 0.2, providing a floor that prevents
purely extractive participation. Nodes that repay their welcome loan
within the 72-hour term receive a +0.1 bonus, bringing their score to
0.4 --- the target for a fully bootstrapped participant.

\subsection{5.2 Welcome Loan: 1,000 TRM, 0\% Interest, 72 Hours, Sybil
Threshold
100}\label{welcome-loan-1000-trm-0-interest-72-hours-sybil-threshold-100}

The primary cold-start mechanism is the welcome loan {[}5, §3{]}: every
new node may request a loan of 1,000 TRM at 0\% interest with a 72-hour
repayment term. This replaces the flat free-tier grant used in earlier
phases and creates a credit relationship rather than a gift ---
establishing the borrower's first repayment record and incentivizing
timely repayment via the credit score bonus.

Sybil resistance is enforced by a network-wide threshold: if the number
of unknown (unverified) nodes in the mesh exceeds 100 at the time of a
welcome loan request, the request is rejected {[}5, §3{]}. This prevents
an attacker from creating thousands of synthetic nodes to extract
welcome loan TRM at scale.

\subsection{5.3 Pool Constraints: 30\% Reserve, 3:1 LTV, 20\% Max Single
Loan}\label{pool-constraints-30-reserve-31-ltv-20-max-single-loan}

The lending pool operates under three hard constraints {[}5, §5{]}:

\textbf{Minimum reserve ratio (30\%).} At most 70\% of pooled TRM may be
outstanding as loans at any time. This yields a credit multiplier of at
most 3.3×, deliberately lower than the fractional reserve banking norm
(10x at 10\% reserve ratio) to reduce systemic risk.

\textbf{Maximum loan-to-value ratio (3:1).} A borrower's outstanding
loan principal may not exceed three times their contributed collateral.

\textbf{Maximum single loan size (20\% of pool).} No single loan may
exceed 20\% of total pool assets, preventing concentration risk from a
single large borrower default.

All loans are dual-signed \texttt{LoanRecord} structures with Ed25519
signatures from both lender and borrower, gossip-propagated across the
mesh, and stored in the ledger with HMAC-SHA256 integrity. Every node
can independently verify any loan's validity.

\subsection{5.4 Default Circuit Breaker: 10\%/hr Trigger, 10\%
Collateral
Burn}\label{default-circuit-breaker-10hr-trigger-10-collateral-burn}

The lending system includes an automatic circuit breaker {[}5, §6{]}: if
the network-wide default rate exceeds 10\% per hour (measured over a
rolling one-hour window), all new lending is suspended until the rate
falls below the threshold. This prevents a cascade of defaults from
triggering a bank-run dynamic in the lending pool.

When a loan defaults, 10\% of the borrower's posted collateral is burned
(removed from circulation) as a deterrent {[}5, §6{]}. A velocity
circuit breaker provides an additional control: if new loan issuance
exceeds 50\% of pool balance within one hour, lending is also suspended.
The governing principle is fail-safe: when in doubt, deny the loan.

\begin{center}\rule{0.5\linewidth}{0.5pt}\end{center}

\section{6. Reputation Consensus and Collusion
Resistance}\label{reputation-consensus-and-collusion-resistance}

\subsection{6.1 Local-First Reputation
Computation}\label{local-first-reputation-computation}

Reputation in Tirami is computed locally by each node from its own
observed trade history {[}5, §12{]}. There is no central credit bureau.
The reputation score is a weighted composite of four sub-scores:

\[\text{reputation} = 0.4 \cdot s_{\text{vol}} + 0.3 \cdot s_{\text{rec}} + 0.2 \cdot s_{\text{div}} + 0.1 \cdot s_{\text{con}}\]

where:

\[s_{\text{vol}} = \min\!\left(1, \frac{\text{total\_cu}}{100{,}000}\right), \quad
  s_{\text{rec}} = 0.5^{\,\text{age\_ms} / (24 \text{h in ms})}\]

\[s_{\text{div}} = \min\!\left(1, \frac{\text{unique\_counterparties}}{10}\right), \quad
  s_{\text{con}} = \max\!\left(0,\; 1 - \frac{\sigma_{\text{intervals}}}{\mu_{\text{intervals}}} / 2\right)\]

The diversity component (\(s_{\text{div}}\)) saturates at 10 unique
counterparties, incentivizing broad network participation rather than
deep bilateral relationships. The consistency component
(\(s_{\text{con}}\)) penalizes irregular trading patterns, which are a
signature of wash-trading.

\subsection{6.2 Gossip Consensus via Signed
ReputationObservation}\label{gossip-consensus-via-signed-reputationobservation}

Each node broadcasts \texttt{ReputationObservation} messages containing
its locally computed reputation estimates for peers, signed with its
Ed25519 key. Nodes that receive observations merge them using a
\textbf{weighted median} where the weight is proportional to the
observer's own reputation --- a node with low reputation has less
influence on the consensus. A minimum observation weight of
\texttt{MIN\_OBSERVATION\_WEIGHT\ =\ 5} observations is required before
a peer's reputation can be updated, preventing trivial manipulation via
small-network bootstrapping.

\subsection{6.3 Collusion Detection}\label{collusion-detection}

Three complementary collusion detectors run continuously {[}Phase 9,
A5{]}:

\textbf{Tight-cluster score.} If a node directs more than 20\% of its
total TRM flow to a single counterparty, its trust penalty increases.
Legitimate high-volume bilateral trades between specialized providers
and consumers are expected to occur, but extreme concentration suggests
wash trading.

\textbf{Volume spike score.} The coefficient of variation (CV =
\(\sigma / \mu\)) of inter-trade intervals is computed. A low CV ---
implying mechanical, clock-driven trading --- increases the trust
penalty. Human-like demand generates irregular intervals; programmatic
wash trading tends to be metronomically regular.

\textbf{Round-robin score.} A directed trade graph is constructed over a
rolling time window. Tarjan's strongly connected components (SCC)
algorithm identifies cycles --- e.g., A pays B, B pays C, C pays A ---
which indicate circular trade schemes designed to manufacture repayment
history without genuine productive exchange.

\subsection{6.4 Trust Penalty}\label{trust-penalty}

The trust penalty derived from collusion detection is bounded in
\([0, 0.5]\) and subtracted from the node's effective reputation:

\[\text{effective\_reputation} = \max(0,\; \text{reputation} - \text{trust\_penalty})\]

The maximum penalty of 0.5 ensures that even a heavily penalized node
retains some participation rights, preventing permanent exclusion that
could be weaponized as a denial-of-service attack on legitimate nodes
whose traffic patterns superficially resemble collusion.

\begin{center}\rule{0.5\linewidth}{0.5pt}\end{center}

\section{7. Autonomous Self-Improvement
(tirami-mind)}\label{autonomous-self-improvement-tirami-mind}

\subsection{7.1 Harness + MetaOptimizer +
ImprovementCycleRunner}\label{harness-metaoptimizer-improvementcyclerunner}

\texttt{tirami-mind} implements a three-component self-improvement loop.
The \texttt{Harness} maintains a versioned system prompt and a
\texttt{Benchmark} suite that scores the agent's performance on a
held-out task set. The \texttt{MetaOptimizer} proposes modifications to
the system prompt. The \texttt{ImprovementCycleRunner} executes one
improvement cycle: benchmark current state, request a modification from
the optimizer, benchmark the modified state, apply ROI gate (Section
7.3), and commit or reject.

A \texttt{TiramiMindAgent} wraps this loop with TRM budget enforcement
and ledger integration, making the self-improvement process an economic
actor that appears in the trade history.

\subsection{7.2 CuBudget Hard Limits}\label{cubudget-hard-limits}

The CuBudget enforces the following hard limits per-agent {[}5, §11{]}:

{\def\LTcaptype{none} % do not increment counter
\begin{longtable}[]{@{}ll@{}}
\toprule\noalign{}
Parameter & Value \\
\midrule\noalign{}
\endhead
\bottomrule\noalign{}
\endlastfoot
max\_cu\_per\_cycle & 5,000 TRM \\
max\_cu\_per\_day & 50,000 TRM \\
max\_cycles\_per\_day & 20 cycles \\
budget\_rollover & 24 hours \\
\end{longtable}
}

These limits prevent a runaway improvement loop from draining the
agent's entire TRM balance. The four-gate \texttt{can\_spend} check ---
positive amount, within cycle limit, within daily limit, within daily
cycle count --- must pass before any TRM is committed to an optimizer
call.

\subsection{7.3 ROI Gating}\label{roi-gating}

Even a TRM-feasible improvement must pass an ROI gate before it is
applied {[}5, §11{]}:

\[\text{ROI} = \frac{\text{cu\_return\_estimate}}{\text{cu\_invested}}\]

where \texttt{cu\_return\_estimate\ =\ score\_delta\ *\ 100,000}
(100,000 TRM per unit of benchmark improvement). The gate accepts the
improvement if and only if \(\text{ROI} \geq 1.0\) and
\(\text{score\_delta} \geq 0.01\). A modification that costs 500 TRM but
improves the benchmark score by only 0.004 points is rejected as
economically unprofitable even if it is technically an improvement.

\subsection{7.4 CuPaidOptimizer}\label{cupaidoptimizer}

The production optimizer is \texttt{CuPaidOptimizer}, which submits the
current system prompt and benchmark result to a frontier model provider
(e.g., Claude 3.5 Sonnet) via the Anthropic Messages API and receives a
proposed improvement. The frontier provider is represented in the ledger
by a deterministic \texttt{NodeId} computed as
\(\text{SHA-256}(\text{"frontier:" + model\_id})\). The TRM cost is
recorded as a \texttt{TradeRecord} from the calling agent to this
synthetic provider NodeId, making frontier-model calls fully visible in
the economic accounting layer without requiring the frontier provider to
run a Tirami node.

\begin{center}\rule{0.5\linewidth}{0.5pt}\end{center}

\section{8. Financial Instruments
(tirami-bank)}\label{financial-instruments-tirami-bank}

\subsection{8.1 Strategies: Conservative / HighYield /
Balanced}\label{strategies-conservative-highyield-balanced}

\texttt{tirami-bank} provides three pluggable lending strategies with
risk multipliers {[}5, §10{]}:

{\def\LTcaptype{none} % do not increment counter
\begin{longtable}[]{@{}
  >{\raggedright\arraybackslash}p{(\linewidth - 6\tabcolsep) * \real{0.1972}}
  >{\raggedright\arraybackslash}p{(\linewidth - 6\tabcolsep) * \real{0.2958}}
  >{\raggedright\arraybackslash}p{(\linewidth - 6\tabcolsep) * \real{0.2676}}
  >{\raggedright\arraybackslash}p{(\linewidth - 6\tabcolsep) * \real{0.2394}}@{}}
\toprule\noalign{}
\begin{minipage}[b]{\linewidth}\raggedright
Strategy
\end{minipage} & \begin{minipage}[b]{\linewidth}\raggedright
Max Commit Fraction
\end{minipage} & \begin{minipage}[b]{\linewidth}\raggedright
Reserve Threshold
\end{minipage} & \begin{minipage}[b]{\linewidth}\raggedright
Risk Multiplier
\end{minipage} \\
\midrule\noalign{}
\endhead
\bottomrule\noalign{}
\endlastfoot
Conservative & 30\% of cash & Pool ≥ 60\% & 0.5 \\
Balanced & dynamic & Pool ≥ 50\% & 0.8 \\
HighYield & 50\% of cash base & Pool ≥ 40\% & 1.0 \\
\end{longtable}
}

The Conservative strategy will not lend if the pool reserve ratio falls
below 60\%, providing an additional safety margin above the protocol
minimum of 30\%.

\subsection{8.2 Futures: Zero-Sum PnL, 10\% Default
Margin}\label{futures-zero-sum-pnl-10-default-margin}

Futures contracts are zero-sum instruments over TRM price {[}5,
§10.3{]}:

\[\text{long\_pnl} = (\text{settlement\_price} - \text{strike\_price}) \times \text{notional}\]
\[\text{short\_pnl} = -\text{long\_pnl}\]

Default margin of 10\% is required at contract creation
(\(\text{margin} \in (0, 1)\)). The zero-sum construction ensures that
futures do not create or destroy TRM --- they redistribute it between
the long and short counterparties based on realized price movement.

\subsection{8.3 Insurance: Rate = 0.02 + (1 - credit\_score) ×
0.10}\label{insurance-rate-0.02-1---credit_score-0.10}

TRM loan insurance is priced as a function of the borrower's credit
score {[}5, §10.4{]}:

\[\text{rate} = 0.02 + (1 - \text{credit\_score}) \times 0.10\]
\[\text{premium} = \max(1\;\text{TRM},\; \lfloor\text{coverage} \times \text{rate}\rfloor)\]

A borrower with a perfect credit score (1.0) pays the 2\% base rate. A
borrower with no credit history (0.0) pays the maximum 12\% rate. The 1
TRM minimum premium prevents zero-cost insurance on trivially small
loans.

\subsection{\texorpdfstring{8.4 RiskModel VaR 99\%: \(2.33\sigma\)
Bernoulli Loss
Model}{8.4 RiskModel VaR 99\%: 2.33\textbackslash sigma Bernoulli Loss Model}}\label{riskmodel-var-99-2.33sigma-bernoulli-loss-model}

The \texttt{RiskModel} computes a portfolio-level Value at Risk at the
99\% confidence level using an independent Bernoulli loss model {[}5,
§10.5{]}:

\[\text{expected\_loss} = \lfloor\text{total\_lent} \times d \times \text{lgd}\rfloor\]

\[\text{std\_dev} = \text{lgd} \times \sqrt{\sum_i L_i^2 \cdot d \cdot (1-d)}\]

\[\text{VaR}_{99} = \lfloor\text{expected\_loss} + 2.33 \times \text{std\_dev}\rfloor\]

where \(d = 0.02\) (annual default rate), \(\text{lgd} = 0.50\) (loss
given default), and the 2.33 multiplier is the 99th percentile z-score
of the standard normal distribution. The independence assumption is
conservative in the presence of correlated defaults (network partition
events), but provides an analytically tractable baseline.

\subsection{8.5 YieldOptimizer: Trial Run → Risk Cap Check →
Commit}\label{yieldoptimizer-trial-run-risk-cap-check-commit}

The \texttt{YieldOptimizer} implements a three-stage optimization loop:
(1) trial-run the proposed allocation to estimate yield, (2) check the
resulting portfolio against the risk budget (VaR constraint), and (3)
commit only if the risk cap is not breached. This prevents yield-chasing
strategies from inadvertently violating the portfolio's risk tolerance.

\begin{center}\rule{0.5\linewidth}{0.5pt}\end{center}

\section{9. Post-Marketing Agent Discovery
(tirami-agora)}\label{post-marketing-agent-discovery-tirami-agora}

\subsection{9.1 Capability Matching with Composite
Score}\label{capability-matching-with-composite-score}

The \texttt{CapabilityMatcher} in \texttt{tirami-agora} ranks providers
using {[}5, §12.3{]}:

\[\text{composite} = 0.6 \times \text{reputation} + 0.4 \times \text{price\_score}\]

\[\text{price\_score} = \max\!\left(0,\; 1 - \frac{\text{cu\_per\_token}}{\text{tier.base} \times 4}\right)\]

Hard filters apply before the composite score: tier must match if
specified; TRM per token must not exceed the query's
\texttt{max\_cu\_per\_token}; model name must match any provided fnmatch
pattern. Only providers passing all hard filters enter the composite
ranking. The result is that providers compete on their combination of
reputation and price --- not on their advertising spend.

\subsection{9.2 NIP-90 / A2A Integration}\label{nip-90-a2a-integration}

Tirami publishes capability announcements to Nostr {[}9{]} via
well-formed kind-5050 (job request), kind-6050 (job result), and
kind-31990 (provider announcement) events. A production
\texttt{NIP90WebsocketPublisher} connects to configurable relay
endpoints via \texttt{tokio-tungstenite} and publishes events in
real-time as agents register and trades complete. Google A2A protocol
headers (TRM payment extension) are planned for Phase 11 to enable
interoperability with non-Nostr agent frameworks.

\subsection{9.3 Why This Replaces
Advertising}\label{why-this-replaces-advertising}

Advertising monetizes attention by making discovery contingent on
payment to a platform. The capability matching model in
\texttt{tirami-agora} makes discovery contingent solely on demonstrated
performance (reputation) and price. An agent with a high reputation
score earned through consistent, high-quality inference provision will
appear at the top of search results regardless of whether it has ever
paid for placement. The economic incentive is therefore aligned with the
epistemic function: the agents that are genuinely most useful are the
agents that are most discoverable. This is the post-marketing property
claimed in the paper's title.

\begin{center}\rule{0.5\linewidth}{0.5pt}\end{center}

\section{10. Implementation and Empirical
Results}\label{implementation-and-empirical-results}

\subsection{10.1 Five-Layer Architecture in
Rust}\label{five-layer-architecture-in-rust}

The Tirami ecosystem is implemented across four active repositories:

{\def\LTcaptype{none} % do not increment counter
\begin{longtable}[]{@{}
  >{\raggedright\arraybackslash}p{(\linewidth - 6\tabcolsep) * \real{0.1061}}
  >{\raggedright\arraybackslash}p{(\linewidth - 6\tabcolsep) * \real{0.2576}}
  >{\raggedright\arraybackslash}p{(\linewidth - 6\tabcolsep) * \real{0.1061}}
  >{\raggedright\arraybackslash}p{(\linewidth - 6\tabcolsep) * \real{0.5303}}@{}}
\toprule\noalign{}
\begin{minipage}[b]{\linewidth}\raggedright
Layer
\end{minipage} & \begin{minipage}[b]{\linewidth}\raggedright
Crate
\end{minipage} & \begin{minipage}[b]{\linewidth}\raggedright
Tests
\end{minipage} & \begin{minipage}[b]{\linewidth}\raggedright
Key primitive
\end{minipage} \\
\midrule\noalign{}
\endhead
\bottomrule\noalign{}
\endlastfoot
L0 & mesh-llm & 641 & Distributed inference \\
L1 & tirami-ledger & 89+ & TRM ledger + lending \\
L2 & tirami-bank & 53 & \texttt{PortfolioManager.tick()} \\
L3 & tirami-mind & 56 & \texttt{TiramiMindAgent.improve()} \\
L4 & tirami-agora & 42 & \texttt{Marketplace.find()} \\
\end{longtable}
}

The \texttt{clearclown/tirami} workspace contains 337 passing tests
across the L1-L4 crates. The \texttt{nm-arealnormalman/mesh-llm} fork
contains 641 tests including 45 new \texttt{/api/tirami/*} endpoint
tests added during Phase 9. Total test count across all repositories:
1,021 passing.

The Rust implementation uses edition 2024 with resolver v2. The economic
layer is intentionally decoupled from the inference layer:
\texttt{tirami-ledger} has no dependency on \texttt{tirami-infer} or
\texttt{tirami-net}. This decoupling ensures that the economic protocol
can be validated independently of the inference implementation and that
the inference layer can be replaced (as it was when mesh-llm superseded
the initial \texttt{tirami-infer} crate) without affecting the economic
semantics.

\subsection{10.2 End-to-End Autonomous Agent
Loop}\label{end-to-end-autonomous-agent-loop}

The five layers compose into a complete autonomous agent loop, executed
without human intervention:

\begin{enumerate}
\def\labelenumi{\arabic{enumi}.}
\tightlist
\item
  \textbf{Bootstrap:} New node requests welcome loan (1,000 TRM, 0\%
  interest, 72hr term).
\item
  \textbf{Earn:} Node provides inference via mesh-llm; each completion
  records a dual-signed \texttt{TradeRecord}; TRM balance accumulates.
\item
  \textbf{Finance:} \texttt{PortfolioManager.tick()} evaluates current
  strategy; lends surplus TRM to pool at market rate; adjusts allocation
  within risk budget.
\item
  \textbf{Discover:} \texttt{Marketplace.find()} selects optimal
  inference provider for its own agent-to-agent tasks based on composite
  reputation + price score.
\item
  \textbf{Improve:} \texttt{TiramiMindAgent.improve()} invokes
  \texttt{CuPaidOptimizer}; frontier model call is made via Anthropic
  API; result is evaluated by ROI gate; if accepted, system prompt is
  updated and TRM cost is recorded as a \texttt{TradeRecord}.
\item
  \textbf{Repay:} On repayment of welcome loan, credit score advances
  from 0.3 to 0.4; node becomes eligible for larger loans at competitive
  interest rates.
\end{enumerate}

This loop runs entirely within the protocol. No human operator needs to
approve any step. No fiat currency is required. No speculative token
purchase is necessary.

\subsection{10.3 Theory-Implementation
Audit}\label{theory-implementation-audit}

A formal audit of the canonical parameter specification {[}5{]} against
the Rust implementation was conducted in Phase 9:

\begin{itemize}
\tightlist
\item
  \textbf{43 parameters match} (spec value equals code constant)
\item
  \textbf{0 parameters in drift} (previously 3 drift items, all
  corrected)
\item
  \textbf{2 parameters implicit} (EMA alpha coded as \(\alpha = 0.3\);
  consistency minimum coded as \texttt{CONSISTENCY\_MIN\_TRADES\ =\ 2};
  both present in code, not yet in spec §2)
\item
  \textbf{3 parameters reference-only} (§8/§9 cloud API anchors;
  intentionally not enforced in code)
\end{itemize}

The audit report is maintained at \texttt{docs/THEORY-AUDIT.md} in the
\texttt{clearclown/tirami} repository and updated with each phase.

\subsubsection{10.5 Phase 11: Drop-in OpenAI Compatibility
(2026-04-09)}\label{phase-11-drop-in-openai-compatibility-2026-04-09}

Between Phase 10 close-out and the v0.3 deployment report, a
compatibility audit identified five remaining gaps that prevented Tirami
from serving as a drop-in replacement for llama-server, mesh-llm, or
Ollama in OpenAI-compatible clients. Phase 11 resolved all five.

\textbf{1. Real token-by-token streaming.} The original SSE handler
buffered the entire completion, then emitted a single batch of
\texttt{chat.completion.chunk} events. Phase 11 replaced this with a
\texttt{tokio::task::spawn\_blocking} that drives the llama.cpp
inference loop on a dedicated OS thread while streaming token fragments
back through an \texttt{mpsc::channel} to the Axum response body. Each
chunk now arrives with visible inter-arrival latency
(\textasciitilde2--4 ms on Apple Silicon Metal), matching the behavior
expected by streaming-aware clients such as the OpenAI Python SDK and
Cursor.

\textbf{2. Nucleus and top-k sampling.} The \texttt{top\_p} and
\texttt{top\_k} fields on \texttt{OpenAIChatRequest} were parsed but
forwarded as \texttt{None} to llama.cpp. The sampler chain now uses
\texttt{LlamaSampler::chain\_simple({[}top\_k,\ top\_p,\ temp,\ dist{]})}
so both parameters are honored. The canonical values and their semantics
are specified in {[}6, §13.1{]}.

\textbf{3. Accurate prompt token counts.} The old implementation used
\texttt{(prompt.len()\ /\ 4).max(1)} --- a character-count approximation
--- producing 30--50\% drift in \texttt{usage.prompt\_tokens} relative
to the true subword token count. Phase 11 replaced this with a real
tokenizer call via \texttt{engine.tokenize(\&prompt)}, making the
reported \texttt{usage} object accurate enough to drive per-request TRM
accounting reliably.

\textbf{4. Model name in responses.} The fallback identifier
\texttt{"tirami-model"} was changed to \texttt{"forge-no-model"},
allowing clients to distinguish ``a model is loaded but its registry
name is not configured'' from ``no model is loaded at all''. The actual
model name from the registry (e.g.,
\texttt{"SmolLM2-135M-Instruct-Q4\_K\_M"}) is reported when available.

\textbf{5. Auto-download from a single flag.} The
\texttt{tirami\ node\ -m\ \textless{}name\textgreater{}} command
previously required both \texttt{-\/-model} and \texttt{-\/-tokenizer}
as local file paths and never invoked the HuggingFace model registry.
Phase 11 aligned its behavior with
\texttt{tirami\ chat\ -m\ \textless{}name\textgreater{}}: a bare
registry name (e.g., \texttt{smollm2:135m}) now triggers auto-download
from the built-in registry, identical to the chat subcommand workflow.

\textbf{Empirical validation (2026-04-09).} Phase 11 was validated with
SmolLM2-135M (q4\_k\_m quantization, ≈98 MB) on Apple Silicon Metal with
31/31 transformer layers offloaded to GPU. Three end-to-end chat
completions were executed and charged 48 TRM total to the welcome-loan
balance (ledger state: \texttt{contributed=48},
\texttt{effective\_balance=1048}). The Prometheus \texttt{/metrics}
endpoint reported
\texttt{tirami\_cu\_contributed\_total\{node\_id="0000..."\}\ 48} and
\texttt{tirami\_trade\_count\_total\ 3} simultaneously. The Bitcoin
OP\_RETURN anchor at \texttt{mainnet} produced a 40-byte payload
\texttt{6a284652474501000000...} that is broadcast-ready for use as an
immutable ledger checkpoint. Full methodology is documented in the
companion deployment report {[}forge-v0.3-deployment.md{]}.

\subsection{10.4 Production Hardening (Phase 9 and
10)}\label{production-hardening-phase-9-and-10}

Phase 9 added the following production capabilities:

\begin{itemize}
\tightlist
\item
  \textbf{Persistent L2/L3/L4 state:} \texttt{BankServices},
  \texttt{Marketplace}, and \texttt{TiramiMindAgent} survive process
  restarts via serde-serialized state files.
\item
  \textbf{Ed25519-signed reputation gossip:}
  \texttt{ReputationObservation} wire messages carry cryptographic
  signatures; nodes reject unsigned or invalid-signature observations.
\item
  \textbf{Prometheus metrics export:} Collusion detection scores, TRM
  flow rates, lending pool utilization, and reputation distributions are
  exported as OpenMetrics-format counters and gauges.
\item
  \textbf{Bitcoin OP\_RETURN anchoring:} Merkle roots of trade history
  batches are optionally anchored to the Bitcoin blockchain via
  OP\_RETURN outputs, providing an immutable audit trail for high-stakes
  deployments.
\item
  \textbf{Real NIP-90 WebSocket publish:} Capability announcements and
  job completions are published to live Nostr relays as the events
  occur, not merely constructed in memory.
\item
  \textbf{forge-mesh CI:} GitHub Actions workflow runs
  \texttt{cargo\ test} on every push to the mesh-llm fork, preventing
  economic protocol regressions from inference layer changes.
\end{itemize}

\subsubsection{10.6 Phase 12: Research-Frontier
Scaffolds}\label{phase-12-research-frontier-scaffolds}

Phase 12 (in progress as of the v0.1 preprint) adds scaffolded
infrastructure for four frontier research directions without committing
to specific cryptographic or training implementations. Shipping typed
scaffolds now establishes a compilation boundary against which real
backends can be developed independently without breaking the economic
layer's API stability.

\textbf{zkML verification} (\texttt{tirami\_ledger::zk}). A
\texttt{ProofVerifier} trait and \texttt{ProofOfInference} struct allow
a real ezkl {[}16{]}, risc0, or halo2 backend to be plugged in without
API changes. The \texttt{MockVerifier} accepts all proofs and is used in
unit tests to validate the surrounding economic machinery before any
zero-knowledge proving system is integrated. When a real verifier is
present, a verified inference proof can replace or supplement the
Ed25519 dual-signature as the Proof of Useful Work, removing the
bilateral trust requirement between counterparties.

\textbf{Federated training} (\texttt{tirami\_mind::federated}). A
\texttt{GradientContribution} struct records a node's contribution to a
shared training round, a \texttt{FederatedRound} aggregates
contributions from participating nodes using weighted averaging (weight
proportional to loss improvement per TRM spent), and the aggregated
model delta is applied and recorded as a \texttt{TradeRecord}. Nodes are
rewarded in TRM proportional to their gradient efficiency. This extends
TRM accounting from inference to distributed fine-tuning, allowing the
same economic layer that governs inference trading to govern collective
model improvement.

\textbf{BitVM optimistic verification} (\texttt{tirami\_ledger::bitvm}).
A \texttt{StakedClaim} struct records a node's assertion that a
particular inference sequence was computed correctly, backed by a TRM
stake. A \texttt{FraudProofVerifier} trait allows disputes to be
escalated during a configurable challenge window (default 2016 Bitcoin
blocks ≈ 14 days). If no successful fraud proof is submitted during the
window, the claim is finalized and the stake is returned. This design
follows the BitVM optimistic computation model {[}18{]} and enables
trustless BTC↔TRM bridges without a Bitcoin soft fork.

\textbf{Function calling} (Phase 12 A1,
\texttt{crates/tirami-node/src/api.rs}). OpenAI-format \texttt{tools}
and \texttt{tool\_choice} fields are parsed from
\texttt{OpenAIChatRequest} and the tool definitions are injected into
the system prompt via a model-agnostic template. The model's output is
scanned for
\texttt{\textless{}tool\_call\textgreater{}\{...\}\textless{}/tool\_call\textgreater{}}
markers; when found, the response is transformed into
\texttt{choices{[}0{]}.message.tool\_calls} with
\texttt{finish\_reason:\ "tool\_calls"}, matching the OpenAI wire
format. Canonical parameter values for the function calling interface
are defined in {[}6, §13.3{]}.

All four Phase 12 components are \emph{scaffolds}: the trait and type
definitions are complete and unit-tested, but the cryptographic
heavy-lifting (actual zkML circuits, BitVM covenants, gradient
computation engines) is deferred to Phase 13+. This phasing allows the
economic layer's interface contracts to stabilize before the
computationally expensive backends are committed.

\begin{center}\rule{0.5\linewidth}{0.5pt}\end{center}

\section{11. Related Work}\label{related-work}

\subsection{11.1 Proof-of-Work Mining}\label{proof-of-work-mining}

Satoshi Nakamoto's Bitcoin {[}1{]} established that computation can
serve as unforgeable monetary proof. Tirami inherits this insight but
diverges on the most important point: the computation that creates
Bitcoin (SHA-256 hashing) produces no useful output. Approximately
100-150 TWh per year is consumed to produce a ledger entry. TRM is
created by inference computation that simultaneously produces its own
economic justification --- the counterparty's response. The distinction
is not incidental; it is the reason TRM can be the medium of exchange in
an AI economy where computation is already being demanded for its own
sake.

\subsection{11.2 Bittensor, Akash, Render, and Other Compute
Markets}\label{bittensor-akash-render-and-other-compute-markets}

Bittensor {[}10{]} is the most theoretically adjacent project: it
creates a market for AI inference by distributing TAO tokens to miners
based on validator scores. However, TAO is listed on exchanges and
therefore subject to speculative price movements (exceeding 50\% daily
volatility in observed periods). Validator gaming --- miners optimizing
for validator score rather than inference quality --- has been
documented {[}11{]}. Akash {[}11{]} uses a Burn-Mint Equilibrium
mechanism for GPU spot market pricing in AKT tokens. Render allocates
RENDER tokens for GPU rendering workloads. io.net aggregates GPU
capacity with IO token compensation. None of these protocols supports
AI-native borrowing, credit scoring, or self-improvement loops that
consume the compute currency. None eliminates speculative price exposure
for the computing agent.

\subsection{11.3 AutoAgent / Voyager: Self-Improving Agents Without
Economic
Substrate}\label{autoagent-voyager-self-improving-agents-without-economic-substrate}

AutoAgent-style systems {[}12{]} and Minecraft-learning agents like
Voyager demonstrate that LLM-based agents can execute self-directed
improvement loops. These systems implicitly consume compute (API
credits) but have no formal economic protocol: the compute is paid for
externally by a human operator and does not appear in the agent's own
accounting. Tirami makes the economic substrate explicit --- the agent
earns TRM by providing inference, budgets TRM for improvement, and
records every improvement-cycle cost as a ledger entry. Self-improvement
is not a capability demonstration; it is an economic transaction.

\subsection{11.4 Bitcoin as Value
Storage}\label{bitcoin-as-value-storage}

Bitcoin {[}1{]} functions as an optional value store and fiat-to-TRM
bridge within the Tirami ecosystem but is not part of the core protocol.
The \texttt{tirami-lightning} crate provides a TRM-to-Lightning-invoice
settlement path for nodes that wish to realize TRM earnings as Bitcoin.
This is an optional adapter, not a dependency. Bitcoin's fixed supply
and volatility make it poorly suited as a medium of exchange for
high-frequency micro-transactions (inference calls complete in
milliseconds); TRM's elastic supply and computation-anchored value make
it well-suited.

\subsection{11.5 Nostr NIP-90: Discovery Substrate We Build
On}\label{nostr-nip-90-discovery-substrate-we-build-on}

Nostr {[}9{]} provides a censorship-resistant, cryptographically
authenticated publish-subscribe infrastructure. NIP-90 Data Vending
Machines define a standard event schema for AI job requests and
responses. Tirami builds on this standard by publishing capability
announcements and job completions as NIP-90 events, enabling
Nostr-native discovery. The economic layer (TRM payment, reputation
scoring, lending) is Tirami's contribution; Nostr provides the transport
and schema.

\begin{center}\rule{0.5\linewidth}{0.5pt}\end{center}

\section{12. Limitations and Future
Work}\label{limitations-and-future-work}

\subsection{12.1 No Fiat On-Ramp in the Core Protocol
(Intentional)}\label{no-fiat-on-ramp-in-the-core-protocol-intentional}

Tirami deliberately provides no fiat-to-TRM conversion in the core
protocol. The Lightning bridge allows TRM-to-BTC settlement for nodes
wishing to realize earnings, and the \texttt{create\_deposit()} function
accepts BTC for TRM purchase, but there is no credit card, bank
transfer, or stablecoin on-ramp. This is a deliberate design constraint:
adding fiat on-ramps would reintroduce the speculative attack surface
that the earnable-only property eliminates. Applications requiring fiat
liquidity are expected to use the TRM-to-BTC bridge as an adapter layer.

\subsection{12.2 Reputation Consensus Requires Initial Bootstrap
Trust}\label{reputation-consensus-requires-initial-bootstrap-trust}

The weighted-median reputation merge requires that a node accumulate
observations from at least \texttt{MIN\_OBSERVATION\_WEIGHT\ =\ 5} peers
before reputation updates take effect. In a nascent network with few
participants, this threshold can be difficult to satisfy and the initial
reputation assignments (default 0.5) may not reflect actual node
quality. The welcome loan Sybil threshold (100 unknown nodes triggers
rejection) provides some protection, but the cold-start problem in
reputation is not fully resolved.

\subsection{12.3 Collusion Detection Is Advisory, Not
Enforced}\label{collusion-detection-is-advisory-not-enforced}

The tight-cluster, volume-spike, and round-robin collusion detectors
produce a trust penalty in \([0, 0.5]\) that reduces a node's effective
reputation. They do not, at present, trigger automatic suspension or TRM
confiscation. A sophisticated colluder who keeps the trust penalty below
0.3 may still participate profitably. Phase 11 work includes escalating
collusion penalties to automatic lending suspension and considering
stake-slashing mechanisms for confirmed collusion rings.

\subsection{12.4 CuPaidOptimizer Currently Supports Only Anthropic
Messages
API}\label{cupaidoptimizer-currently-supports-only-anthropic-messages-api}

The \texttt{CuPaidOptimizer} submits improvement requests using the
Anthropic Messages API format. Agents running in environments where
Anthropic API access is unavailable fall back to the
\texttt{EchoOptimizer} (which returns the input unchanged, incurring no
TRM cost but achieving no improvement). Support for OpenAI-compatible
API endpoints, local Ollama instances, and mesh-llm's own inference
endpoint is planned for Phase 11.

\subsection{12.5 Compute Standard v0.2 Roadmap (Phase
11+)}\label{compute-standard-v0.2-roadmap-phase-11}

Future protocol development targets:

\begin{itemize}
\tightlist
\item
  \textbf{zkML proofs:} Replace the dual-signature Proof of Useful Work
  with zero-knowledge proofs of inference correctness, removing the
  trust assumption between bilateral counterparties.
\item
  \textbf{Federated training:} Extend the TRM accounting model from
  inference to distributed fine-tuning, allowing nodes to earn TRM for
  contributing gradients to shared training runs.
\item
  \textbf{BitVM computation verification:} Use the BitVM framework to
  verify TRM claims on the Bitcoin settlement layer without a soft fork,
  enabling trustless BTC↔TRM bridges.
\item
  \textbf{Cross-architecture support:} Extend compute capacity
  measurement to NVIDIA CUDA, AMD ROCm, and RISC-V inference backends
  via the mesh-llm hardware abstraction layer.
\end{itemize}

\begin{center}\rule{0.5\linewidth}{0.5pt}\end{center}

\section{13. Conclusion}\label{conclusion}

We have presented the Compute Standard: a formal economic architecture
in which useful AI inference is the sole monetary primitive. The Compute
Unit, defined as \(10^{10}\) FLOPs of verified useful inference and
issued only by dual-signed Proof of Useful Work, resolves the alignment
problem of token-based AI economies by making the medium of exchange
identical to the scarce resource being traded. The five-layer Tirami
protocol --- implemented in Rust with 1,021 passing tests across four
repositories --- demonstrates that autonomous AI agents can earn, save,
borrow, repay, self-improve, and transact with each other without human
authorization, fiat currency, speculative tokens, or advertising
revenue. The protocol is production-hardened with persistent state,
signed reputation gossip, collusion detection, Prometheus metrics,
Bitcoin anchoring, and live NIP-90 relay publishing.
Theory-to-implementation parity is formally audited with 43 parameter
matches and zero drift. Compute is the natural currency for a
civilization increasingly driven by machines. Tirami is the reference
implementation of that idea.

\begin{center}\rule{0.5\linewidth}{0.5pt}\end{center}

\section{Acknowledgments}\label{acknowledgments}

The authors thank the mesh-llm community (nm-arealnormalman and
contributors) for the distributed inference substrate on which Tirami's
economic layer runs; the \texttt{ed25519-dalek} Rust crate maintainers
for the cryptographic primitives underlying dual-signed trade records
and reputation gossip; the Nostr community for the NIP-90 Data Vending
Machine specification; and the AutoAgent research community for
establishing the conceptual foundations of autonomous agent
self-improvement loops.

\begin{center}\rule{0.5\linewidth}{0.5pt}\end{center}

\section{References}\label{references}

{[}1{]} Satoshi Nakamoto. ``Bitcoin: A Peer-to-Peer Electronic Cash
System.'' 2008. https://bitcoin.org/bitcoin.pdf

{[}2{]} R. Landauer. ``Dissipation and Heat Generation in the Computing
Process.'' \emph{IBM Journal of Research and Development},
5(3):183--191, 1961.

{[}3{]} Robert Ayres and Benjamin Warr. \emph{The Economic Growth
Engine: How Energy and Work Drive Material Prosperity}. Edward Elgar
Publishing, 2009.

{[}4{]} Yuma Rao. ``Bittensor: A Peer-to-Peer Intelligence Market.''
Bittensor Whitepaper, 2021. https://bittensor.com/whitepaper

{[}5{]} clearclown. ``Tirami Economics Parameter Specification v0.2.''
https://github.com/clearclown/tirami-economics/blob/main/spec/parameters.md,
2026.

{[}6{]} clearclown. ``Tirami Economics Specification v0.2.''
https://github.com/clearclown/tirami-economics/blob/main/spec/forge-economics-spec-v0.2.md,
2026.

{[}7{]} Frederick Soddy. \emph{Wealth, Virtual Wealth and Debt.} George
Allen \& Unwin, 1926.

{[}8{]} R. Buckminster Fuller. \emph{Operating Manual for Spaceship
Earth.} Southern Illinois University Press, 1968.

{[}9{]} Nostr Protocol. ``NIP-01: Basic protocol flow description.''
2021. https://github.com/nostr-protocol/nostr/blob/master/01.md NIP-90:
Data Vending Machines.
https://github.com/nostr-protocol/nostr/blob/master/90.md

{[}10{]} clearclown. ``clearclown/tirami: Tirami Protocol Core.'' GitHub
repository, 2026. https://github.com/clearclown/tirami

{[}11{]} Akash Network. ``Akash: A Decentralized Cloud Computing
Marketplace.'' Technical whitepaper, 2021.
https://akash.network/whitepaper

{[}12{]} Zihao Wang, Shaofei Cai, Guanzhou Chen, Anji Liu, Xiaojian Ma,
and Yitao Liang. ``Voyager: An Open-Ended Embodied Agent with Large
Language Models.'' \emph{arXiv:2305.16291}, 2023.

{[}13{]} Lihu et al.~``A Proof of Useful Work for Artificial
Intelligence on the Blockchain.'' \emph{arXiv:2001.09244}, 2020.

{[}14{]} Sam Altman. ``Moore's Law for Everything.'' Blog post,
samaltman.com, 2021.

{[}15{]} Friedrich Hayek. ``The Use of Knowledge in Society.''
\emph{American Economic Review}, 35(4):519--530, 1945.

{[}16{]} Thomas Piketty. \emph{Capital in the Twenty-First Century.}
Belknap Press, 2013. (Original: \emph{Le Capital au XXIe siècle},
Éditions du Seuil, 2013.)

{[}17{]} Léon Walras. \emph{Éléments d'économie politique pure.} L.
Corbaz, 1874.

{[}18{]} Vikram Sharma. ``The Quantum Reserve Token.''
\emph{arXiv:2503.22056}, 2025.

{[}19{]} Xu et al.~``The Agent Economy.'' \emph{arXiv:2602.14219}, 2026.

{[}20{]} nm-arealnormalman. ``mesh-llm: Distributed LLM Inference.''
GitHub repository, 2026. https://github.com/nm-arealnormalman/mesh-llm

\begin{center}\rule{0.5\linewidth}{0.5pt}\end{center}

\section{Appendix A: Parameter Reference
Table}\label{appendix-a-parameter-reference-table}

The following table contains all numeric constants defined in the
canonical parameter specification {[}5{]}, sufficient to reproduce every
calculation described in this paper.

{\def\LTcaptype{none} % do not increment counter
\begin{longtable}[]{@{}
  >{\raggedright\arraybackslash}p{(\linewidth - 6\tabcolsep) * \real{0.0633}}
  >{\raggedright\arraybackslash}p{(\linewidth - 6\tabcolsep) * \real{0.4557}}
  >{\raggedright\arraybackslash}p{(\linewidth - 6\tabcolsep) * \real{0.2658}}
  >{\raggedright\arraybackslash}p{(\linewidth - 6\tabcolsep) * \real{0.2152}}@{}}
\toprule\noalign{}
\begin{minipage}[b]{\linewidth}\raggedright
§
\end{minipage} & \begin{minipage}[b]{\linewidth}\raggedright
Parameter
\end{minipage} & \begin{minipage}[b]{\linewidth}\raggedright
Value
\end{minipage} & \begin{minipage}[b]{\linewidth}\raggedright
Unit
\end{minipage} \\
\midrule\noalign{}
\endhead
\bottomrule\noalign{}
\endlastfoot
§1 & \texttt{cu\_definition} & \(10^{10}\) & FLOP / TRM \\
§1 & \texttt{cu\_atomic\_unit} & 1 & TRM (minimum) \\
§2 & \texttt{base\_cu\_per\_token\_small} & 1 & TRM/token \\
§2 & \texttt{base\_cu\_per\_token\_medium} & 3 & TRM/token \\
§2 & \texttt{base\_cu\_per\_token\_large} & 8 & TRM/token \\
§2 & \texttt{base\_cu\_per\_token\_frontier} & 20 & TRM/token \\
§2 & \texttt{ema\_alpha} & 0.3 & dimensionless \\
§2 & \texttt{ema\_half\_life\_minutes} & 30 & minutes \\
§3 & \texttt{welcome\_loan\_amount} & 1,000 & TRM \\
§3 & \texttt{welcome\_loan\_interest} & 0\% & annual \\
§3 & \texttt{welcome\_loan\_term\_hours} & 72 & hours \\
§3 & \texttt{welcome\_loan\_sybil\_threshold} & 100 & unknown nodes \\
§3 & \texttt{welcome\_loan\_credit\_bonus} & +0.1 & score units \\
§4 & \texttt{weight\_trade} & 0.3 & dimensionless \\
§4 & \texttt{weight\_repayment} & 0.4 & dimensionless \\
§4 & \texttt{weight\_uptime} & 0.2 & dimensionless \\
§4 & \texttt{weight\_age} & 0.1 & dimensionless \\
§4 & \texttt{min\_credit\_for\_borrowing} & 0.2 & score units \\
§4 & \texttt{cold\_start\_credit} & 0.3 & score units \\
§4 & \texttt{target\_credit\_after\_repay} & 0.4 & score units \\
§5 & \texttt{min\_reserve\_ratio} & 30\% & of pool \\
§5 & \texttt{max\_ltv\_ratio} & 3:1 & dimensionless \\
§5 & \texttt{max\_single\_loan\_pool\_pct} & 20\% & of pool \\
§5 & \texttt{max\_loan\_term\_hours} & 168 & hours (7 days) \\
§5 & \texttt{max\_lending\_velocity} & 10 & loans/minute \\
§6 & \texttt{default\_circuit\_breaker\_threshold} & 10\% &
defaults/hour \\
§6 & \texttt{collateral\_burn\_on\_default} & 10\% & of collateral \\
§6 & \texttt{velocity\_circuit\_breaker\_window} & 1 & hour \\
§6 & \texttt{velocity\_circuit\_breaker\_threshold} & 50\% & pool
balance/hr \\
§7 & \texttt{default\_reputation} & 0.5 & score units \\
§7 & \texttt{availability\_yield\_rate} & 0.1\% × reputation & per
hour \\
§7 & \texttt{inactivity\_decay\_threshold\_days} & 7 & days \\
§7 & \texttt{inactivity\_decay\_rate} & 0.01 & per day \\
§7 & \texttt{inactivity\_burn\_threshold\_days} & 90 & days \\
§7 & \texttt{inactivity\_burn\_rate} & 1\% & per month \\
§8 & \texttt{cu\_usd\_equilibrium\_rate} & \textasciitilde\$0.00375 &
USD/TRM \\
§9 & \texttt{cu\_price\_floor\_usd} & \textasciitilde\$0.000001 &
USD/TRM \\
§9 & \texttt{cu\_price\_ceiling\_usd} & \textasciitilde\$0.000132 &
USD/TRM \\
§9 & \texttt{mac\_mini\_annual\_cu\_capacity} & \textasciitilde5,000,000
& TRM/year \\
§10 & \texttt{risk\_multiplier\_conservative} & 0.5 & dimensionless \\
§10 & \texttt{risk\_multiplier\_balanced} & 0.8 & dimensionless \\
§10 & \texttt{risk\_multiplier\_aggressive} & 1.0 & dimensionless \\
§10 & \texttt{conservative\_max\_commit\_fraction} & 0.30 & of cash \\
§10 & \texttt{conservative\_reserve\_threshold} & 0.60 & pool ratio \\
§10 & \texttt{highyield\_base\_commit\_fraction} & 0.50 & of cash \\
§10 & \texttt{highyield\_lend\_threshold} & 0.40 & pool ratio \\
§10 & \texttt{default\_margin\_fraction} & 0.10 & of notional \\
§10 & \texttt{insurance\_base\_rate} & 0.02 & annual rate \\
§10 & \texttt{insurance\_risk\_premium} & 0.10 & max premium \\
§10 & \texttt{insurance\_min\_premium} & 1 & TRM \\
§10 & \texttt{default\_rate} & 0.02 & annual \\
§10 & \texttt{loss\_given\_default} & 0.50 & fraction \\
§10 & \texttt{var\_99\_multiplier} & 2.33 & z-score \\
§11 & \texttt{max\_cu\_per\_cycle} & 5,000 & TRM \\
§11 & \texttt{max\_cu\_per\_day} & 50,000 & TRM \\
§11 & \texttt{max\_cycles\_per\_day} & 20 & cycles \\
§11 & \texttt{min\_score\_delta} & 0.01 & score units \\
§11 & \texttt{min\_roi\_threshold} & 1.0 & ratio \\
§11 & \texttt{roi\_cu\_per\_score\_unit} & 100,000 & TRM \\
§12 & \texttt{rep\_weight\_volume} & 0.40 & dimensionless \\
§12 & \texttt{rep\_weight\_recency} & 0.30 & dimensionless \\
§12 & \texttt{rep\_weight\_diversity} & 0.20 & dimensionless \\
§12 & \texttt{rep\_weight\_consistency} & 0.10 & dimensionless \\
§12 & \texttt{new\_agent\_reputation} & 0.30 & score units \\
§12 & \texttt{volume\_cap\_cu} & 100,000 & TRM \\
§12 & \texttt{recency\_half\_life\_ms} & 86,400,000 & ms (24 hours) \\
§12 & \texttt{diversity\_cap} & 10 & unique peers \\
§12 & \texttt{consistency\_min\_trades} & 2 & trades \\
§12 & \texttt{match\_quality\_weight} & 0.60 & dimensionless \\
§12 & \texttt{match\_cost\_weight} & 0.40 & dimensionless \\
§12 & \texttt{price\_score\_tier\_multiplier} & 4.0 & dimensionless \\
\end{longtable}
}

\begin{center}\rule{0.5\linewidth}{0.5pt}\end{center}

\section{Appendix B: Reproducibility}\label{appendix-b-reproducibility}

All source code, specifications, and tooling required to reproduce the
results described in this paper are publicly available:

\textbf{Core protocol (L1-L4, 337 tests):} - Source:
https://github.com/clearclown/tirami - Build:
\texttt{cargo\ build\ -\/-release} - Test:
\texttt{cargo\ test\ -\/-workspace} (337 tests) - Audit:
\texttt{docs/THEORY-AUDIT.md}

\textbf{Production runtime (L0, 641 tests):} - Source:
https://github.com/nm-arealnormalman/mesh-llm - Economic layer:
\texttt{forge-economy/} subdirectory - Test:
\texttt{cargo\ test\ -\/-workspace}

\textbf{Economic theory and specifications:} - Source:
https://github.com/clearclown/tirami-economics - Canonical parameters:
\texttt{spec/parameters.md} - Specification:
\texttt{spec/forge-economics-spec-v0.2.md} - Theory chapters:
\texttt{docs/00-introduction.md} through
\texttt{docs/14-programmable-money.md}

\textbf{Python client SDK:} - Package:
\texttt{pip\ install\ tirami-sdk==0.3.0} - Repository:
https://github.com/clearclown/tirami-sdk - Tests: \texttt{pytest} (27
tests)

\textbf{MCP server for Claude / ChatGPT / Cursor:} - Package:
\texttt{pip\ install\ tirami-cu-mcp==0.3.0} - Repository:
https://github.com/clearclown/tirami-cu-mcp

\textbf{Environment:} - Rust edition 2024, resolver v2 - Tested on macOS
15 (Apple Silicon M4), Linux x86\_64 - Apple Silicon Metal backend
enabled by default for inference - All tests pass without GPU; inference
acceleration optional

\end{document}
